% from ega0-18-fr.tex
\oldpage[0\textsubscript{IV-1}]{65}

les classes de $A$-extensions équivalentes de $B$ par $L$ qui sont des
$A$-\emph{algèbres commutatives} (ou, ce qui revient au même, des
$A$-\emph{anneaux commutatifs}) forment un \emph{sous-groupe} de
$\operatorname{Exal}_A(B,L)$, noté $\operatorname{Exalcom}_A(B,L)$. Si
$\mathbf K''$ est la sous-catégorie pleine de $\mathbf K'$ formée des triplets
$(A,B,L)$ où $A$ et $B$ sont commutatifs et $L$ un $B$-module,
\[
  (A,B,L)\longmapsto\operatorname{Exalcom}_A(B,L)
\]
est encore un \emph{foncteur covariant} de $\mathbf K''$ dans $\mathbf{Ab}$. On
peut aussi dans (18.3.7) remplacer partout $\operatorname{Exan}$ par
$\operatorname{Exalcom}$. Enfin, si dans (18.3.8), on suppose que $B$ soit un
anneau \emph{commutatif} et que $L$ soit un $C$-\emph{module}, le même
raisonnement donne un isomorphisme canonique
\begin{equation}
\label{0.18.4.2.1-fr}
\tag{18.4.2.1}
  \operatorname{Hom}_C(\mathfrak J/\mathfrak J^2,L)
  \xrightarrow{\sim}\operatorname{Exalcom}_B(C,L)
\end{equation}
où le premier membre est le groupe des homomorphismes de $C$-\emph{module}.
\end{env}

\begin{env}[18.4.3]
\label{0.18.4.3-fr}
Soit $A$ un anneau commutatif. Un cas important d'extensions de
$A$-algèbres est formé des $A$-algèbres $E$, extensions d'une $A$-algèbre $B$
par un $B$-bimodule $L$ telles que $E$, en tant que $A$-\emph{module}, soit une
extension \emph{triviale} du $A$-module $B$ par le $A$-module $L$~; autrement
dit, la suite exacte de $A$-\emph{modules}
$0\to L\to E\to B\to0$ est \emph{scindée}. On dit alors que $E$ est une
\emph{extension de Hochschild} de $B$ par $L$. Ce sera toujours le cas lorsque
$B$ est un $A$-module \emph{projectif}, et en particulier lorsque $A$ est un
\emph{corps} commutatif. En tant que $A$-module, on peut identifier $E$ à
$B\times L$, la multiplication dans $E$ étant donnée par
$(x,s)(y,t)=(xy,xt+sy+f(x,y))$ avec $f(x,y)\in L$. Si l'on écrit que cette
multiplication définit sur $B\times L$ une structure de $A$-algèbre, on trouve
(M, XIV, 2) que $f$ doit être une application $A$-\emph{bilinéaire} de
$B\times B$ dans $L$, telle que
\begin{equation}
\label{0.18.4.3.1-fr}
\tag{18.4.3.1}
  f(xy,z)+f(x,y)z=xf(y,z)+f(x,yz),
\end{equation}
autrement dit, $f$ est un $2$-\emph{cocycle} sur $B$ à valeurs dans $L$, au
sens de Hochschild~; pour que l'extension $E$ soit $A$-triviale, il faut et il
suffit que l'on ait
\begin{equation}
\label{0.18.4.3.2-fr}
\tag{18.4.3.2}
  f(x,y)=xg(y)-g(xy)+g(x)y,
\end{equation}
où $g$ est une application $A$-\emph{linéaire} de $B$ dans $L$, autrement dit
$f$ doit être un $2$-\emph{cobord} au sens de Hochschild. On en déduit aussitôt
que les classes d'extensions de Hochschild de $B$ par $L$ forment un
sous-groupe de $\operatorname{Exal}_A(B,L)$, isomorphe au \emph{groupe de
cohomologie} de Hochschild $\mathrm H_A^2(B,L)$.

Si $B$ est une $A$-algèbre \emph{commutative}, et $L$ un $B$-module, les
extensions de Hochschild \emph{commutatives} de $B$ par $L$ correspondent aux
$2$-cocycles \emph{symétriques}, c'est-à-dire tels que $f(y,x)=f(x,y)$. Les
classes d'extensions de Hochschild commutatives de $B$ par $L$ forment donc un
sous-groupe de $\operatorname{Exalcom}_A(B,L)$, isomorphe au sous-groupe de
$\mathrm H_A^2(B,L)$ image du groupe des $2$-cocycles symétriques, que nous
noterons $\mathrm H_A^2(B,L)^s$.
\end{env}

\begin{env}[18.4.4]
\label{0.18.4.4-fr}
Bornons-nous au cas où $A$ et $B$ sont \emph{commutatifs} et $L$ un
$B$-module, et rappelons dans ce cas la définition équivalente des groupes de
Hochschild $\mathrm H_A^i(B,L)$ (M, IX, 6). On considère un complexe
$\mathbf P_\bullet=(\mathbf P_n)_{n\geqslant0}$ de $B$-modules, où
$\mathbf P_n=B^{\otimes(n+1)}=B\otimes_A B\otimes_A\cdots\otimes_A B$
($n+1$ fois), la structure de $B$-module étant définie

